% Bloom levels: remember, understand, apply, analyze, evaluate, create.
\begin{problema}\label{prob:3f1ec07}
State, without calculation, Pearson's chi-square goodness-of-fit statistic, and the general rule for determining its degrees of freedom when the theoretical model is fully specified in advance, with no parameters estimated from the sample.
\end{problema}

\begin{problema}\label{prob:d4859df}
A data scientist designs a $\chi^2$ goodness-of-fit test for a discrete categorization model with four classes ($k=4$, $\nu=3$) and $n=80$ observations, obtaining expected frequencies $E_1=45.2$, $E_2=30.5$, $E_3=3.2$, and $E_4=1.1$. Explain why directly computing Pearson's statistic violates the asymptotic regularity conditions known as Cochran's rule, and propose the standard correction procedure, indicating how many degrees of freedom the adjusted test would have.
\end{problema}

\begin{problema}\label{prob:da2eb15}
A geneticist investigates whether ABO blood groups in a cohort follow national reference proportions: $45\%$ A, $40\%$ B, $10\%$ AB, and $5\%$ O. In a sample of $n=200$ people, the observed counts are A=95, B=75, AB=18, O=12. Perform the goodness-of-fit test at level $\alpha=0.05$, using $\chi^2_{0.05,3}=7.815$.
\end{problema}

\begin{problema}\label{prob:d6f0154}
An engineer records the daily number of server interruptions during $n=100$ days: $0$ interruptions on 36 days, $1$ on 40 days, $2$ on 18 days, and $\ge3$ on 6 days, treating the last category as the point value $k=3$. To check whether interruptions follow a Poisson model: (a) compute $\hat\lambda$ by maximum likelihood; (b) compute the expected frequencies under Poisson($\hat\lambda$); (c) compute $\chi^2$ and justify why the correct degrees of freedom are $\nu=k-1-m$ with $m=1$ estimated parameter. Is the model accepted at $\alpha=0.05$?
\end{problema}

\begin{problema}\label{prob:f622374}
In $2\times2$ tables with small or medium samples, Pearson's uncorrected $\chi^2$ statistic tends to produce systematically lower $p$-values than Fisher's exact test, increasing Type I error risk. Evaluate why this bias occurs by explaining the difference between the discrete, stepwise nature of the exact hypergeometric distribution and the continuous nature of the $\chi^2_1$ reference distribution, and explain how Yates' continuity correction, using $|ad-bc|-N/2$ instead of $|ad-bc|$, corrects this bias.
\end{problema}

\begin{problema}\label{prob:369a6a1}
Design an original example, with context, categories, and hypothetical observed frequencies, where a $\chi^2$ goodness-of-fit test is applied. Use a context different from dice, blood types, or servers. State $H_0$ and $H_a$, compute the expected frequencies, and compute the $\chi^2$ statistic.
\end{problema}

\begin{solproblema}[prob:3f1ec07]
\[
	\chi^2=\sum_{i=1}^k\frac{(O_i-E_i)^2}{E_i}.
\]
When the model is fully specified and no parameters are estimated from the sample, the degrees of freedom are $\nu=k-1$, where $k$ is the number of categories.
\end{solproblema}

\begin{solproblema}[prob:d4859df]
The continuous $\chi^2$ approximation to the true discrete distribution of
\[
	Q=\sum\frac{(O_i-E_i)^2}{E_i}
\]
is reliable only when no expected frequency is below $1$ and no more than $20\%$ of cells have expected frequency below $5$. Here $E_3=3.2$ and $E_4=1.1$, so $2$ of $4$ cells, or $50\%$, violate the $E_i<5$ rule. The standard correction is to combine adjacent low-frequency categories. Combining classes 3 and 4 gives $E_{3+4}=4.3$, resulting in $k'=3$ effective categories and $\nu=k'-1=2$ degrees of freedom.
\end{solproblema}

\begin{solproblema}[prob:da2eb15]
The expected counts are
\[
	E_A=90,\qquad E_B=80,\qquad E_{AB}=20,\qquad E_O=10.
\]
Then
\[
	\chi^2=\frac{(95-90)^2}{90}+\frac{(75-80)^2}{80}
	+\frac{(18-20)^2}{20}+\frac{(12-10)^2}{10}
	=1.190.
\]
Since $1.190<7.815$, do not reject $H_0$: the sample is consistent with the national reference proportions.
\end{solproblema}

\begin{solproblema}[prob:d6f0154]
\begin{enumerate}[label=(\alph*)]
	\item
	\[
		\hat\lambda=\bar X=\frac{0(36)+1(40)+2(18)+3(6)}{100}=0.94.
	\]
	\item Under Poisson($0.94$), the expected frequencies are approximately
	\[
		E_0=39.06,\quad E_1=36.72,\quad E_2=17.26,\quad E_{\ge3}=6.96.
	\]
	\item
	\[
		\chi^2\approx
		\frac{(36-39.06)^2}{39.06}
		+\frac{(40-36.72)^2}{36.72}
		+\frac{(18-17.26)^2}{17.26}
		+\frac{(6-6.96)^2}{6.96}
		\approx0.697.
	\]
	Because one parameter, $\hat\lambda$, was estimated from the sample, the degrees of freedom are $\nu=4-1-1=2$. Since $0.697\ll\chi^2_{0.05,2}=5.991$, do not reject the Poisson model.
\end{enumerate}
\end{solproblema}

\begin{solproblema}[prob:f622374]
Fisher's exact reference distribution for a finite $2\times2$ table is hypergeometric, so it is discrete and stepwise. The $\chi^2_1$ reference is continuous. Approximating a stepwise tail by a continuous curve can make the right-tail area, the $p$-value, too small, producing too many false positives. Yates' correction reduces the numerator by replacing $|ad-bc|$ with $|ad-bc|-N/2$, which lowers the statistic and increases the approximate $p$-value toward the exact Fisher value.
\end{solproblema}

\begin{solproblema}[prob:369a6a1]
Example: an online store checks whether product returns follow historical category proportions: Clothing $40\%$, Electronics $30\%$, Home $20\%$, Other $10\%$. In $n=150$ returns, the observed counts are Clothing=68, Electronics=40, Home=25, Other=17. The hypotheses are that the proportions are $40/30/20/10\%$ versus at least one differs. Expected counts are $E=[60,45,30,15]$. Thus
\[
	\chi^2=\frac{(68-60)^2}{60}+\frac{(40-45)^2}{45}
	+\frac{(25-30)^2}{30}+\frac{(17-15)^2}{15}
	\approx2.723.
\]
With $\chi^2_{0.05,3}=7.815$, do not reject $H_0$.
\end{solproblema}
